\documentclass[12pt,a4paper]{article}
\usepackage{ugmath}
\usepackage{float}
\usepackage{placeins}
\usepackage{array}
\usepackage[labelfont=bf,labelsep=space]{caption}
\newcommand{\paren}[1]{\left( #1 \right)}
\newcommand{\alumno}{Ricardo León Martínez}
\newcommand{\materia}{Algebra Lineal I}
\newcommand{\profesor}{Claudia Reynoso Alcántara}
\newcommand{\tarea}{Tarea 12}
\newcommand{\fecha}{21/5/2026}


\begin{document}

    \begin{center}
        {\large\textbf{UNIVERSIDAD DE GUANAJUATO}}\\[0.3cm]
        {\normalsize\textbf{DIVISIÓN DE CIENCIAS NATURALES Y EXACTAS}}\\
        {\normalsize\textbf{CAMPUS GUANAJUATO}}\\[1cm]

        {\Large\textbf{\tarea\ (\materia)}}\\[1cm]
    \end{center}

    \noindent
    \textbf{Nombre:} \alumno 
    \hfill 
    \textbf{Fecha:} \fecha 
    \hfill 
    \textbf{Calificación:} \rule{3cm}{0.4pt}

    \vspace{0.3cm}

    \begin{mdframed}[style=mdpinkbox,frametitle={Ejercicio 1}]
        Sean $A,B\in M_{n\times n}(F)$. Demuestra que
        \begin{enumerate}
            \item[(a)] $(A^{t})^{t}=A$,
            \item[(b)] $(A+B)^{t}=A^{t}+B^{t}$,
            \item[(c)] $(AB)^{t}=B^{t}A^{t}$.  
        \end{enumerate}
    \end{mdframed}

    \begin{proof}
        \begin{enumerate}
            \item[(a)] $(A^{t})^{t}=A$\\
            Para cualesquiera $i,j=\{1,\dots,n\}$,
            \begin{equation*}
                ((A^{t})^{t})_{ij}=(A^{t})_{ji}=a_{ij}.
            \end{equation*}
            Por lo tanto,
            \begin{equation*}
                ((A^{t})^{t})_{ij}=a_{ij}=A_{ij}
            \end{equation*}
            para todo $i,j$. Luego,
            \begin{equation*}
                (A^{t})^{t}=A.
            \end{equation*}
            \item[(b)] $(A+B)^{t}=A^{t}+B^{t}$\\
            Sea $i,j\in\{1,\dots,n\}$. Entonces,
            \begin{equation*}
                ((A+B)^{t})_{ij}=(A+B)_{ji}.
            \end{equation*}
            Como la suma de matrices se define entrada a entrada,
            \begin{equation*}
                (A+B)_{ji}=a_{ji}+b_{ji}.
            \end{equation*}
            Así,
            \begin{equation*}
                ((A+B)^{t})_{ij}=(A^{t})_{ij}+(B^{t})_{ij}=(A^{t}+B^{t})_{ij}.
            \end{equation*}
            Conluimos que
            \begin{equation*}
                (A+B)^{t}=A^{t}+B^{t}.
            \end{equation*}
            \item[(c)] $(AB)^{t}=B^{t}A^{t}$\\
            Tomemos $i,j\in\{1,\dots,n\}$. Entonces
            \begin{equation*}
                ((AB)^{t})_{ij}=(AB)_{ij}.
            \end{equation*}
            Por definición del producto matricial,
            \begin{equation*}
                (AB)_{ji}=\sum_{k=1}^{n}a_{jk}b_{ki}.
            \end{equation*}
            Ahora observamos que
            \begin{equation*}
                a_{jk}=(A^{t})_{jk},\qquad b_{ki}=(B^{t})_{ik}.
            \end{equation*}
            Sustituyendo,
            \begin{equation*}
                (AB)_{ji}=\sum_{k=1}^{n}(B^{t})_{ik}(A^{t})_{kj}.
            \end{equation*}
            Pero, por definición de producto matricial.
            \begin{equation*}
                \sum_{k=1}^{n}(B^{t})_{ik}(A^{t})_{kj}=(B^{t}A^{t})_{ij}.
            \end{equation*}
            Por consiguiente,
            \begin{equation*}
                ((AB)^{t})_{ij}=(B^{t}A^{t})_{ij}.
            \end{equation*}
            Por lo tanto,
            \begin{equation*}
                (AB)^{t}=B^{t}A^{t}.
            \end{equation*}
        \end{enumerate}
    \end{proof}

    \begin{mdframed}[style=mdpinkbox,frametitle={Ejercicio 2}]
        Sea $D:M_{2\times2}(F)\to F$ una función alternante y 2-lineal. Demuestra que para toda
        matriz $A$:
        \begin{equation*}
            D(A)=\det(A)D(I_{2}).
        \end{equation*}
    \end{mdframed}

    \begin{proof}
        Sea
        \begin{equation*}
            A=
            \begin{pmatrix}
                a & b\\
                c & d
            \end{pmatrix}
            \in M_{2\times 2}(F).
        \end{equation*}
        Denotemos por
        \begin{equation*}
            v_{1}=
            \begin{pmatrix}
                a\\
                c
            \end{pmatrix},
            v_{2}=
            \begin{pmatrix}
                b\\
                d
            \end{pmatrix}
        \end{equation*}
        las columnas de $A$. Entonces
        \begin{equation*}
            A(v_{1},v_{2}).
        \end{equation*}
        Como $D$ es 2-lineal, es lineal en cada columna. Además,
        \begin{equation*}
            v_{1}=ae_{1}+ce_{2},\qquad v_{2}=be_{1}+de_{2},
        \end{equation*}
        donde
        \begin{equation*}
            e_{1}=
            \begin{pmatrix}
                1\\
                0
            \end{pmatrix},
            \qquad
            e_{2}=
            \begin{pmatrix}
                0\\
                1
            \end{pmatrix}.
        \end{equation*}
        Por binilinealidad,
        \begin{align*}
            D(A)&=D(ae_{1}+ce_{2},be_{1}+de_{2})\\
            &=abD(e_{1},e_{1})+adD(e_{1},e_{2})+cbD(e_{2},e_{1})+cdD(e_{2},e_{2}).
        \end{align*}
        Ahora usamos que $D$ es alternante. Por definición,
        \begin{equation*}
            D(e_{1},e_{1})=0,\qquad D(e_{2},e_{2})=0.
        \end{equation*}
        Ademas, una función alternante satisface
        \begin{equation*}
            D(u,v)=-D(v,u)
        \end{equation*}
        para todos $u,v$. Por tanto,
        \begin{equation*}
            D(e_{2},e_{1})=-D(e_{1},e_{2}).
        \end{equation*}
        Sustituyendo,
        \begin{align*}
            D(A)&=adD(e_{1},e_{2})+cb(-D(e_{1},e_{2}))\\
            &=(ad-bc)D(e_{1},e_{2}).
        \end{align*}
        Observemos ahora que
        \begin{equation*}
            (e_{1},e_{2})=
            \begin{pmatrix}
                1 & 0\\
                0 & 1
            \end{pmatrix}
            =I_{2}.
        \end{equation*}
        Luego,
        \begin{equation*}
            D(e_{1},e_{2})=D(I_{2}).
        \end{equation*}
        Como
        \begin{equation*}
            \det(A)=ad-bc,
        \end{equation*}
        concluimos que
        \begin{equation*}
            D(A)=\det(A)D(I_{n}).
        \end{equation*}
    \end{proof}

    \begin{mdframed}[style=mdpinkbox,frametitle={Ejercicio 3}]
        Sea $A\in M_{n\times n}(F)$, demuestra que $\det(A)=\det(A^{t})$
    \end{mdframed}

    \begin{proof}
        Sea $A=(a_{ij})\in M_{n\times n}(F)$. Por definición del determinante,
        \begin{equation*}
            \det(A)=\sum_{\sigma\in S_{n}}\sgn(\sigma)\prod_{i=1}^{n}(A^{t})_{i,\sigma(i)}.
        \end{equation*}
        donde $S_{n}$ es el grupo simetrico de todas las permutaciones de $\{1,\dots, n\}$.
        Ahora consideremos la matriz transpuesta $A^{t}$. Como
        \begin{equation*}
            (A^{t})_{ij}=a_{ji},
        \end{equation*}
        se tiene
        \begin{equation*}
            \det(A^{t})=\sum_{\sigma\in S_{n}}\sgn(\sigma)\prod_{i=1}^{n}(A^{t})_{i,\sigma(i)}.
        \end{equation*}
        Sustituyendo la definición de transpuesta,
        \begin{equation*}
            \det(A^{t})=\sum_{\sigma\in S_{n}}\sgn(\sigma)\prod_{i=1}^{n}(a)_{\sigma(i),i}.
        \end{equation*}
        Ahora hacemos el cambio de variable
        \begin{equation*}
            \tau=\sigma^{-1}.
        \end{equation*}
        Como el inverso de una permutación tiene el mismo signo,
        \begin{equation*}
            \sgn(\tau)=\sgn(\sigma),
        \end{equation*}
        y además, al recorrer $\sigma\in S_{n}$, también recorremos todas las $\tau\in S_{n}$.
        Entonces,
        \begin{equation*}
            \det(A^{t})=\sum_{\tau\in S_{n}}\sgn(\tau)\prod_{i=1}^{n}(a)_{i,\tau(i)}.
        \end{equation*}
        Pero esta ultima expresión es exactamente $\det(A)$. Por lo tanto
        \begin{equation*}
            \det(A^{t})=\det(A).
        \end{equation*}
    \end{proof}

    \begin{mdframed}[style=mdpinkbox,frametitle={Ejercicio 4}]
        Considera la matriz de Vandermonde
        \begin{equation*}
            A=
            \begin{pmatrix}
                1 & a & a^{2}\\
                1 & b & b^{2}\\
                1 & c & c^{2}
            \end{pmatrix},
        \end{equation*}
        sobre $F$. Demuestra que es invertible si y solo si $a\neq b$ o $a\neq c$ o $b\neq c$.
    \end{mdframed}

    \begin{proof}
        Sabemos que una matriz cuadrada es invertible si y solo si su determinante es distinto
        de cero. Por tanto, basta calcular $\det(A)$. Desarrollamos el determinante
        \begin{equation*}
            \det(A)=\det
            \begin{pmatrix}
                1 & a & a^{2}\\
                1 & b & b^{2}\\
                1 & c & c^{2}
            \end{pmatrix}
        \end{equation*}
        Hacemos las operaciones elementales
        \begin{equation*}
            R_{2}\rightarrow R_{2}-R_{1},\qquad R_{3}\rightarrow R_{3}-R_{1},
        \end{equation*}
        Entonces,
        \begin{equation*}
            \det(A)=\det
            \begin{pmatrix}
                1 & a & a^{2}\\
                0 & b-a & b^{2}-a^{2}\\
                0 & c & c^{2}-a^{2}
            \end{pmatrix}
        \end{equation*}
        Factorizamos diferencias de cuadrados
        \begin{equation*}
            b^{2}-a^{2}=(b-a)(b+a),\qquad c^{2}-a^{2}=(c-a)(c+a).
        \end{equation*}
        Así,
        \begin{equation*}
            \det(A)=\det
            \begin{pmatrix}
                1 & a & a^{2}\\
                0 & b-a & (b-a)(b+a)\\
                0 & c-a & (c-a)(c+a)
            \end{pmatrix}
        \end{equation*}
        Sacando los factores comunes de las filas 2 y 3,
        \begin{equation*}
            \det(A)=(b-a)(c-a)\det
            \begin{pmatrix}
                1 & a & a^{2}\\
                0 & 1 & b+a\\
                0 & 1 & c+a
            \end{pmatrix}
        \end{equation*}
        Ahora desarrollamos por cofactores respecto a la primera columna
        \begin{equation*}
            \det(A)=(b-a)(c-a)\det
            \begin{pmatrix}
                1 & b+a\\
                1 & c+a
            \end{pmatrix}.
        \end{equation*}
        Calculando el determinante $2\times2$,
        \begin{equation*}
            \det(A)=(b-a)(c-a)((c+a)-(b+a))
        \end{equation*}
        Simplificando,
        \begin{equation*}
            \det(A)=(b-a)(c-a)(c-b).
        \end{equation*}
        Por tanto
        \begin{equation*}
            \det(A)\neq0\ssi(b-a)(c-a)(c-b)\neq0.
        \end{equation*}
        Esto ocurre exactamente cuando
        \begin{equation*}
            a\neq b,\qquad a\neq c,\qquad b\neq c.
        \end{equation*}
        Así, la matriz $A$ es invertible si y solo si $a\neq b$, $b\neq c$ y $c\neq a$.
    \end{proof}

    \begin{mdframed}[style=mdpinkbox,frametitle={Ejercicio 5}]
        Sea $n\in\N$, $\sigma\in S_{n}$, demuestra que:
        \begin{equation*}
            \sgn(\sigma)=\det
            \begin{pmatrix}
                \epsilon_{\sigma(1)}\\
                \vdots\\
                \epsilon_{\sigma(n)}
            \end{pmatrix}
            =\prod_{1\leq i\lt j\leq n}\frac{\sigma(i)-\sigma(j)}{i-j},
        \end{equation*}
        donde $\epsilon_{i}$ es la fila con todas las entradas cero, excepto el lugar $i$
        que tiene un 1. 
    \end{mdframed}

    \begin{proof}
        Sea $\sigma\in S_{n}$. Consideremos la matriz
        \begin{equation*}
            P_{\sigma}=
            \begin{pmatrix}
                \varepsilon_{\sigma(1)}\\
                \vdots\\
                \varepsilon_{\sigma(n)}
            \end{pmatrix},
        \end{equation*}
        donde $\e_{i}$ es el vector fila canónico con un 1 en la posición $i$ y ceros en las demás
        entradas. Demostremos que
        \begin{equation*}
            \sgn(\sigma)=\det(P_{\sigma}).
        \end{equation*}
        Por definición del determinante
        \begin{equation*}
            \det(P_{\sigma})=\sum_{\tau\in S_{n}}\sgn(\tau)\prod_{i=1}^{n}(P_{\sigma})_{i,\tau(i)}.
        \end{equation*}
        Pero la fila $i$ de $P_{\sigma}$ tiene un único 1, situado en la columna $\sigma(i)$.
        Entonces,
        \begin{equation*}
            (P_{\sigma})_{i,\tau(i)}=
            \begin{cases}
                1,\quad\tau(i)=\sigma(i)\\
                0,\quad\text{en cualquier otro caso}
            \end{cases}
        \end{equation*}
        Por tanto
        \begin{equation*}
            \prod_{i=1}^{n}(P_{\sigma})_{i,\tau(i)}=0.
        \end{equation*}
        para toda $\tau\neq\sigma$, y únicamente el término $\tau=\sigma$. Así,
        \begin{equation*}
            \det(P_{\sigma})=\sgn(\sigma).
        \end{equation*}
        Concluimos que
        \begin{equation*}
            \sgn(\sigma)=\det
            \begin{pmatrix}
                \e_{\sigma(1)}\\
                \vdots\\
                \e_{\sigma(n)}
            \end{pmatrix}.
        \end{equation*}
        Demostremos ahora que
        \begin{equation*}
            \sgn(\sigma)=\prod_{1\leq i\lt j\leq n}\frac{\sigma(i)-\sigma(j)}{i-j}.
        \end{equation*}
        Definamos
        \begin{equation*}
            P=\prod_{1\leq i\lt j\leq n}(i-j).
        \end{equation*}
        Aplicando la permutación $\sigma$ a los números $1\dots, n$, obtenemos
        \begin{equation*}
            \prod_{1\leq i\lt j\leq n}(\sigma(i)-\sigma(j)).
        \end{equation*}
        Observemos que cada vez que intercambiamos dos elementos, el producto cambia de signo,
        porque exactamente uno de los factores cambia de signo. Ahora, toda permutación puede
        escribirse como el producto de transposiciones. Si $\sigma$ se escribe como el producto
        de $k$ transposiciones, entonces
        \begin{equation*}
            Q=(-1)^{k}P.
        \end{equation*}
        Pero por definición,
        \begin{equation*}
            \sgn(\sigma)=(-1)^{k}.
        \end{equation*}
        Así,
        \begin{equation*}
            Q=\sgn(\sigma)P.
        \end{equation*}
        Es decir,
        \begin{equation*}
            \prod_{1\leq i\lt j\leq n}(\sigma(i)-\sigma(j))=\sgn(\sigma)\prod_{1\leq i\lt j\leq n}(i-j).
        \end{equation*}
        Dividiendo entre el producto del denominador,
        \begin{equation*}
            \sgn(\sigma)=\prod_{1\leq i\lt j\leq n}\frac{\sigma(i)-\sigma(j)}{i-j}.
        \end{equation*}
    \end{proof}

\end{document}